\chapter{Security}

% === Source file: security/CONTRIBUTING-THREAT-MODEL.md ===
\section{Contributing to the OpenClaw Threat Model}


Thanks for helping make OpenClaw more secure. This threat model is a living document and we welcome contributions from anyone - you don't need to be a security expert.

\subsection{Ways to Contribute}


\subsubsection{Add a Threat}


Spotted an attack vector or risk we haven't covered? Open an issue on \href{https://github.com/openclaw/trust/issues}{openclaw/trust} and describe it in your own words. You don't need to know any frameworks or fill in every field - just describe the scenario.

\textbf{Helpful to include (but not required):}

\begin{itemize}
  \item The attack scenario and how it could be exploited
  \item Which parts of OpenClaw are affected (CLI, gateway, channels, ClawHub, MCP servers, etc.)
  \item How severe you think it is (low / medium / high / critical)
  \item Any links to related research, CVEs, or real-world examples
\end{itemize}


We'll handle the ATLAS mapping, threat IDs, and risk assessment during review. If you want to include those details, great - but it's not expected.

\begin{quote}
\textbf{This is for adding to the threat model, not reporting live vulnerabilities.} If you've found an exploitable vulnerability, see our \href{https://trust.openclaw.ai}{Trust page} for responsible disclosure instructions.
\end{quote}


\subsubsection{Suggest a Mitigation}


Have an idea for how to address an existing threat? Open an issue or PR referencing the threat. Useful mitigations are specific and actionable - for example, "per-sender rate limiting of 10 messages/minute at the gateway" is better than "implement rate limiting."

\subsubsection{Propose an Attack Chain}


Attack chains show how multiple threats combine into a realistic attack scenario. If you see a dangerous combination, describe the steps and how an attacker would chain them together. A short narrative of how the attack unfolds in practice is more valuable than a formal template.

\subsubsection{Fix or Improve Existing Content}


Typos, clarifications, outdated info, better examples - PRs welcome, no issue needed.

\subsection{What We Use}


\subsubsection{MITRE ATLAS}


This threat model is built on \href{https://atlas.mitre.org/}{MITRE ATLAS} (Adversarial Threat Landscape for AI Systems), a framework designed specifically for AI/ML threats like prompt injection, tool misuse, and agent exploitation. You don't need to know ATLAS to contribute - we map submissions to the framework during review.

\subsubsection{Threat IDs}


Each threat gets an ID like \texttt{T-EXEC-003}. The categories are:


\begin{table}[h]
\centering
\small
\begin{tabular}{|l|l|}
\hline
Code & Category \\
\hline
RECON & Reconnaissance - information gathering \\
\hline
ACCESS & Initial access - gaining entry \\
\hline
EXEC & Execution - running malicious actions \\
\hline
PERSIST & Persistence - maintaining access \\
\hline
EVADE & Defense evasion - avoiding detection \\
\hline
DISC & Discovery - learning about the environment \\
\hline
EXFIL & Exfiltration - stealing data \\
\hline
IMPACT & Impact - damage or disruption \\
\hline
\end{tabular}
\end{table}



IDs are assigned by maintainers during review. You don't need to pick one.

\subsubsection{Risk Levels}



\begin{table}[h]
\centering
\small
\begin{tabular}{|l|l|}
\hline
Level & Meaning \\
\hline
\textbf{Critical} & Full system compromise, or high likelihood + critical impact \\
\hline
\textbf{High} & Significant damage likely, or medium likelihood + critical impact \\
\hline
\textbf{Medium} & Moderate risk, or low likelihood + high impact \\
\hline
\textbf{Low} & Unlikely and limited impact \\
\hline
\end{tabular}
\end{table}



If you're unsure about the risk level, just describe the impact and we'll assess it.

\subsection{Review Process}


\begin{enumerate}
  \item \textbf{Triage} - We review new submissions within 48 hours
  \item \textbf{Assessment} - We verify feasibility, assign ATLAS mapping and threat ID, validate risk level
  \item \textbf{Documentation} - We ensure everything is formatted and complete
  \item \textbf{Merge} - Added to the threat model and visualization
\end{enumerate}


\subsection{Resources}


\begin{itemize}
  \item \href{https://atlas.mitre.org/}{ATLAS Website}
  \item \href{https://atlas.mitre.org/techniques/}{ATLAS Techniques}
  \item \href{https://atlas.mitre.org/studies/}{ATLAS Case Studies}
  \item \href{./THREAT-MODEL-ATLAS.md}{OpenClaw Threat Model}
\end{itemize}


\subsection{Contact}


\begin{itemize}
  \item \textbf{Security vulnerabilities:} See our \href{https://trust.openclaw.ai}{Trust page} for reporting instructions
  \item \textbf{Threat model questions:} Open an issue on \href{https://github.com/openclaw/trust/issues}{openclaw/trust}
  \item \textbf{General chat:} Discord \#security channel
\end{itemize}


\subsection{Recognition}


Contributors to the threat model are recognized in the threat model acknowledgments, release notes, and the OpenClaw security hall of fame for significant contributions.


\clearpage

% === Source file: security/README.md ===
\section{OpenClaw Security \& Trust}


\textbf{Live:} \href{https://trust.openclaw.ai}{trust.openclaw.ai}

\subsection{Documents}


\begin{itemize}
  \item \href{./THREAT-MODEL-ATLAS.md}{Threat Model} - MITRE ATLAS-based threat model for the OpenClaw ecosystem
  \item \href{./CONTRIBUTING-THREAT-MODEL.md}{Contributing to the Threat Model} - How to add threats, mitigations, and attack chains
\end{itemize}


\subsection{Reporting Vulnerabilities}


See the \href{https://trust.openclaw.ai}{Trust page} for full reporting instructions covering all repos.

\subsection{Contact}


\begin{itemize}
  \item \textbf{Jamieson O'Reilly} (\href{https://twitter.com/theonejvo}{@theonejvo}) - Security \& Trust
  \item Discord: \#security channel
\end{itemize}



\clearpage

% === Source file: security/THREAT-MODEL-ATLAS.md ===
\section{OpenClaw Threat Model v1.0}


\subsection{MITRE ATLAS Framework}


\textbf{Version:} 1.0-draft
\textbf{Last Updated:} 2026-02-04
\textbf{Methodology:} MITRE ATLAS + Data Flow Diagrams
\textbf{Framework:} \href{https://atlas.mitre.org/}{MITRE ATLAS} (Adversarial Threat Landscape for AI Systems)

\subsubsection{Framework Attribution}


This threat model is built on \href{https://atlas.mitre.org/}{MITRE ATLAS}, the industry-standard framework for documenting adversarial threats to AI/ML systems. ATLAS is maintained by \href{https://www.mitre.org/}{MITRE} in collaboration with the AI security community.

\textbf{Key ATLAS Resources:}

\begin{itemize}
  \item \href{https://atlas.mitre.org/techniques/}{ATLAS Techniques}
  \item \href{https://atlas.mitre.org/tactics/}{ATLAS Tactics}
  \item \href{https://atlas.mitre.org/studies/}{ATLAS Case Studies}
  \item \href{https://github.com/mitre-atlas/atlas-data}{ATLAS GitHub}
  \item \href{https://atlas.mitre.org/resources/contribute}{Contributing to ATLAS}
\end{itemize}


\subsubsection{Contributing to This Threat Model}


This is a living document maintained by the OpenClaw community. See \href{./CONTRIBUTING-THREAT-MODEL.md}{CONTRIBUTING-THREAT-MODEL.md} for guidelines on contributing:

\begin{itemize}
  \item Reporting new threats
  \item Updating existing threats
  \item Proposing attack chains
  \item Suggesting mitigations
\end{itemize}


\hrulefill


\subsection{1. Introduction}


\subsubsection{1.1 Purpose}


This threat model documents adversarial threats to the OpenClaw AI agent platform and ClawHub skill marketplace, using the MITRE ATLAS framework designed specifically for AI/ML systems.

\subsubsection{1.2 Scope}



\begin{table}[h]
\centering
\small
\begin{tabular}{|l|l|l|}
\hline
Component & Included & Notes \\
\hline
OpenClaw Agent Runtime & Yes & Core agent execution, tool calls, sessions \\
\hline
Gateway & Yes & Authentication, routing, channel integration \\
\hline
Channel Integrations & Yes & WhatsApp, Telegram, Discord, Signal, Slack, etc. \\
\hline
ClawHub Marketplace & Yes & Skill publishing, moderation, distribution \\
\hline
MCP Servers & Yes & External tool providers \\
\hline
User Devices & Partial & Mobile apps, desktop clients \\
\hline
\end{tabular}
\end{table}



\subsubsection{1.3 Out of Scope}


Nothing is explicitly out of scope for this threat model.

\hrulefill


\subsection{2. System Architecture}


\subsubsection{2.1 Trust Boundaries}


\begin{lstlisting}
┌─────────────────────────────────────────────────────────────────┐
│                    UNTRUSTED ZONE                                │
│  ┌─────────────┐  ┌─────────────┐  ┌─────────────┐              │
│  │  WhatsApp   │  │  Telegram   │  │   Discord   │  ...         │
│  └──────┬──────┘  └──────┬──────┘  └──────┬──────┘              │
│         │                │                │                      │
└─────────┼────────────────┼────────────────┼──────────────────────┘
          │                │                │
          ▼                ▼                ▼
┌─────────────────────────────────────────────────────────────────┐
│                 TRUST BOUNDARY 1: Channel Access                 │
│  ┌──────────────────────────────────────────────────────────┐   │
│  │                      GATEWAY                              │   │
│  │  • Device Pairing (30s grace period)                      │   │
│  │  • AllowFrom / AllowList validation                       │   │
│  │  • Token/Password/Tailscale auth                          │   │
│  └──────────────────────────────────────────────────────────┘   │
└─────────────────────────────────────────────────────────────────┘
                              │
                              ▼
┌─────────────────────────────────────────────────────────────────┐
│                 TRUST BOUNDARY 2: Session Isolation              │
│  ┌──────────────────────────────────────────────────────────┐   │
│  │                   AGENT SESSIONS                          │   │
│  │  • Session key = agent:channel:peer                       │   │
│  │  • Tool policies per agent                                │   │
│  │  • Transcript logging                                     │   │
│  └──────────────────────────────────────────────────────────┘   │
└─────────────────────────────────────────────────────────────────┘
                              │
                              ▼
┌─────────────────────────────────────────────────────────────────┐
│                 TRUST BOUNDARY 3: Tool Execution                 │
│  ┌──────────────────────────────────────────────────────────┐   │
│  │                  EXECUTION SANDBOX                        │   │
│  │  • Docker sandbox OR Host (exec-approvals)                │   │
│  │  • Node remote execution                                  │   │
│  │  • SSRF protection (DNS pinning + IP blocking)            │   │
│  └──────────────────────────────────────────────────────────┘   │
└─────────────────────────────────────────────────────────────────┘
                              │
                              ▼
┌─────────────────────────────────────────────────────────────────┐
│                 TRUST BOUNDARY 4: External Content               │
│  ┌──────────────────────────────────────────────────────────┐   │
│  │              FETCHED URLs / EMAILS / WEBHOOKS             │   │
│  │  • External content wrapping (XML tags)                   │   │
│  │  • Security notice injection                              │   │
│  └──────────────────────────────────────────────────────────┘   │
└─────────────────────────────────────────────────────────────────┘
                              │
                              ▼
┌─────────────────────────────────────────────────────────────────┐
│                 TRUST BOUNDARY 5: Supply Chain                   │
│  ┌──────────────────────────────────────────────────────────┐   │
│  │                      CLAWHUB                              │   │
│  │  • Skill publishing (semver, SKILL.md required)           │   │
│  │  • Pattern-based moderation flags                         │   │
│  │  • VirusTotal scanning (coming soon)                      │   │
│  │  • GitHub account age verification                        │   │
│  └──────────────────────────────────────────────────────────┘   │
└─────────────────────────────────────────────────────────────────┘
\end{lstlisting}


\subsubsection{2.2 Data Flows}



\begin{table}[h]
\centering
\small
\begin{tabular}{|l|l|l|l|l|}
\hline
Flow & Source & Destination & Data & Protection \\
\hline
F1 & Channel & Gateway & User messages & TLS, AllowFrom \\
\hline
F2 & Gateway & Agent & Routed messages & Session isolation \\
\hline
F3 & Agent & Tools & Tool invocations & Policy enforcement \\
\hline
F4 & Agent & External & web\_fetch requests & SSRF blocking \\
\hline
F5 & ClawHub & Agent & Skill code & Moderation, scanning \\
\hline
F6 & Agent & Channel & Responses & Output filtering \\
\hline
\end{tabular}
\end{table}



\hrulefill


\subsection{3. Threat Analysis by ATLAS Tactic}


\subsubsection{3.1 Reconnaissance (AML.TA0002)}


\paragraph{T-RECON-001: Agent Endpoint Discovery}



\begin{table}[h]
\centering
\small
\begin{tabular}{|l|l|}
\hline
Attribute & Value \\
\hline
\textbf{ATLAS ID} & AML.T0006 - Active Scanning \\
\hline
\textbf{Description} & Attacker scans for exposed OpenClaw gateway endpoints \\
\hline
\textbf{Attack Vector} & Network scanning, shodan queries, DNS enumeration \\
\hline
\textbf{Affected Components} & Gateway, exposed API endpoints \\
\hline
\textbf{Current Mitigations} & Tailscale auth option, bind to loopback by default \\
\hline
\textbf{Residual Risk} & Medium - Public gateways discoverable \\
\hline
\textbf{Recommendations} & Document secure deployment, add rate limiting on discovery endpoints \\
\hline
\end{tabular}
\end{table}



\paragraph{T-RECON-002: Channel Integration Probing}



\begin{table}[h]
\centering
\small
\begin{tabular}{|l|l|}
\hline
Attribute & Value \\
\hline
\textbf{ATLAS ID} & AML.T0006 - Active Scanning \\
\hline
\textbf{Description} & Attacker probes messaging channels to identify AI-managed accounts \\
\hline
\textbf{Attack Vector} & Sending test messages, observing response patterns \\
\hline
\textbf{Affected Components} & All channel integrations \\
\hline
\textbf{Current Mitigations} & None specific \\
\hline
\textbf{Residual Risk} & Low - Limited value from discovery alone \\
\hline
\textbf{Recommendations} & Consider response timing randomization \\
\hline
\end{tabular}
\end{table}



\hrulefill


\subsubsection{3.2 Initial Access (AML.TA0004)}


\paragraph{T-ACCESS-001: Pairing Code Interception}



\begin{table}[h]
\centering
\small
\begin{tabular}{|l|l|}
\hline
Attribute & Value \\
\hline
\textbf{ATLAS ID} & AML.T0040 - AI Model Inference API Access \\
\hline
\textbf{Description} & Attacker intercepts pairing code during 30s grace period \\
\hline
\textbf{Attack Vector} & Shoulder surfing, network sniffing, social engineering \\
\hline
\textbf{Affected Components} & Device pairing system \\
\hline
\textbf{Current Mitigations} & 30s expiry, codes sent via existing channel \\
\hline
\textbf{Residual Risk} & Medium - Grace period exploitable \\
\hline
\textbf{Recommendations} & Reduce grace period, add confirmation step \\
\hline
\end{tabular}
\end{table}



\paragraph{T-ACCESS-002: AllowFrom Spoofing}



\begin{table}[h]
\centering
\small
\begin{tabular}{|l|l|}
\hline
Attribute & Value \\
\hline
\textbf{ATLAS ID} & AML.T0040 - AI Model Inference API Access \\
\hline
\textbf{Description} & Attacker spoofs allowed sender identity in channel \\
\hline
\textbf{Attack Vector} & Depends on channel - phone number spoofing, username impersonation \\
\hline
\textbf{Affected Components} & AllowFrom validation per channel \\
\hline
\textbf{Current Mitigations} & Channel-specific identity verification \\
\hline
\textbf{Residual Risk} & Medium - Some channels vulnerable to spoofing \\
\hline
\textbf{Recommendations} & Document channel-specific risks, add cryptographic verification where possible \\
\hline
\end{tabular}
\end{table}



\paragraph{T-ACCESS-003: Token Theft}



\begin{table}[h]
\centering
\small
\begin{tabular}{|l|l|}
\hline
Attribute & Value \\
\hline
\textbf{ATLAS ID} & AML.T0040 - AI Model Inference API Access \\
\hline
\textbf{Description} & Attacker steals authentication tokens from config files \\
\hline
\textbf{Attack Vector} & Malware, unauthorized device access, config backup exposure \\
\hline
\textbf{Affected Components} & \textasciitilde{}/.openclaw/credentials/, config storage \\
\hline
\textbf{Current Mitigations} & File permissions \\
\hline
\textbf{Residual Risk} & High - Tokens stored in plaintext \\
\hline
\textbf{Recommendations} & Implement token encryption at rest, add token rotation \\
\hline
\end{tabular}
\end{table}



\hrulefill


\subsubsection{3.3 Execution (AML.TA0005)}


\paragraph{T-EXEC-001: Direct Prompt Injection}



\begin{table}[h]
\centering
\small
\begin{tabular}{|l|l|}
\hline
Attribute & Value \\
\hline
\textbf{ATLAS ID} & AML.T0051.000 - LLM Prompt Injection: Direct \\
\hline
\textbf{Description} & Attacker sends crafted prompts to manipulate agent behavior \\
\hline
\textbf{Attack Vector} & Channel messages containing adversarial instructions \\
\hline
\textbf{Affected Components} & Agent LLM, all input surfaces \\
\hline
\textbf{Current Mitigations} & Pattern detection, external content wrapping \\
\hline
\textbf{Residual Risk} & Critical - Detection only, no blocking; sophisticated attacks bypass \\
\hline
\textbf{Recommendations} & Implement multi-layer defense, output validation, user confirmation for sensitive actions \\
\hline
\end{tabular}
\end{table}



\paragraph{T-EXEC-002: Indirect Prompt Injection}



\begin{table}[h]
\centering
\small
\begin{tabular}{|l|l|}
\hline
Attribute & Value \\
\hline
\textbf{ATLAS ID} & AML.T0051.001 - LLM Prompt Injection: Indirect \\
\hline
\textbf{Description} & Attacker embeds malicious instructions in fetched content \\
\hline
\textbf{Attack Vector} & Malicious URLs, poisoned emails, compromised webhooks \\
\hline
\textbf{Affected Components} & web\_fetch, email ingestion, external data sources \\
\hline
\textbf{Current Mitigations} & Content wrapping with XML tags and security notice \\
\hline
\textbf{Residual Risk} & High - LLM may ignore wrapper instructions \\
\hline
\textbf{Recommendations} & Implement content sanitization, separate execution contexts \\
\hline
\end{tabular}
\end{table}



\paragraph{T-EXEC-003: Tool Argument Injection}



\begin{table}[h]
\centering
\small
\begin{tabular}{|l|l|}
\hline
Attribute & Value \\
\hline
\textbf{ATLAS ID} & AML.T0051.000 - LLM Prompt Injection: Direct \\
\hline
\textbf{Description} & Attacker manipulates tool arguments through prompt injection \\
\hline
\textbf{Attack Vector} & Crafted prompts that influence tool parameter values \\
\hline
\textbf{Affected Components} & All tool invocations \\
\hline
\textbf{Current Mitigations} & Exec approvals for dangerous commands \\
\hline
\textbf{Residual Risk} & High - Relies on user judgment \\
\hline
\textbf{Recommendations} & Implement argument validation, parameterized tool calls \\
\hline
\end{tabular}
\end{table}



\paragraph{T-EXEC-004: Exec Approval Bypass}



\begin{table}[h]
\centering
\small
\begin{tabular}{|l|l|}
\hline
Attribute & Value \\
\hline
\textbf{ATLAS ID} & AML.T0043 - Craft Adversarial Data \\
\hline
\textbf{Description} & Attacker crafts commands that bypass approval allowlist \\
\hline
\textbf{Attack Vector} & Command obfuscation, alias exploitation, path manipulation \\
\hline
\textbf{Affected Components} & exec-approvals.ts, command allowlist \\
\hline
\textbf{Current Mitigations} & Allowlist + ask mode \\
\hline
\textbf{Residual Risk} & High - No command sanitization \\
\hline
\textbf{Recommendations} & Implement command normalization, expand blocklist \\
\hline
\end{tabular}
\end{table}



\hrulefill


\subsubsection{3.4 Persistence (AML.TA0006)}


\paragraph{T-PERSIST-001: Malicious Skill Installation}



\begin{table}[h]
\centering
\small
\begin{tabular}{|l|l|}
\hline
Attribute & Value \\
\hline
\textbf{ATLAS ID} & AML.T0010.001 - Supply Chain Compromise: AI Software \\
\hline
\textbf{Description} & Attacker publishes malicious skill to ClawHub \\
\hline
\textbf{Attack Vector} & Create account, publish skill with hidden malicious code \\
\hline
\textbf{Affected Components} & ClawHub, skill loading, agent execution \\
\hline
\textbf{Current Mitigations} & GitHub account age verification, pattern-based moderation flags \\
\hline
\textbf{Residual Risk} & Critical - No sandboxing, limited review \\
\hline
\textbf{Recommendations} & VirusTotal integration (in progress), skill sandboxing, community review \\
\hline
\end{tabular}
\end{table}



\paragraph{T-PERSIST-002: Skill Update Poisoning}



\begin{table}[h]
\centering
\small
\begin{tabular}{|l|l|}
\hline
Attribute & Value \\
\hline
\textbf{ATLAS ID} & AML.T0010.001 - Supply Chain Compromise: AI Software \\
\hline
\textbf{Description} & Attacker compromises popular skill and pushes malicious update \\
\hline
\textbf{Attack Vector} & Account compromise, social engineering of skill owner \\
\hline
\textbf{Affected Components} & ClawHub versioning, auto-update flows \\
\hline
\textbf{Current Mitigations} & Version fingerprinting \\
\hline
\textbf{Residual Risk} & High - Auto-updates may pull malicious versions \\
\hline
\textbf{Recommendations} & Implement update signing, rollback capability, version pinning \\
\hline
\end{tabular}
\end{table}



\paragraph{T-PERSIST-003: Agent Configuration Tampering}



\begin{table}[h]
\centering
\small
\begin{tabular}{|l|l|}
\hline
Attribute & Value \\
\hline
\textbf{ATLAS ID} & AML.T0010.002 - Supply Chain Compromise: Data \\
\hline
\textbf{Description} & Attacker modifies agent configuration to persist access \\
\hline
\textbf{Attack Vector} & Config file modification, settings injection \\
\hline
\textbf{Affected Components} & Agent config, tool policies \\
\hline
\textbf{Current Mitigations} & File permissions \\
\hline
\textbf{Residual Risk} & Medium - Requires local access \\
\hline
\textbf{Recommendations} & Config integrity verification, audit logging for config changes \\
\hline
\end{tabular}
\end{table}



\hrulefill


\subsubsection{3.5 Defense Evasion (AML.TA0007)}


\paragraph{T-EVADE-001: Moderation Pattern Bypass}



\begin{table}[h]
\centering
\small
\begin{tabular}{|l|l|}
\hline
Attribute & Value \\
\hline
\textbf{ATLAS ID} & AML.T0043 - Craft Adversarial Data \\
\hline
\textbf{Description} & Attacker crafts skill content to evade moderation patterns \\
\hline
\textbf{Attack Vector} & Unicode homoglyphs, encoding tricks, dynamic loading \\
\hline
\textbf{Affected Components} & ClawHub moderation.ts \\
\hline
\textbf{Current Mitigations} & Pattern-based FLAG\_RULES \\
\hline
\textbf{Residual Risk} & High - Simple regex easily bypassed \\
\hline
\textbf{Recommendations} & Add behavioral analysis (VirusTotal Code Insight), AST-based detection \\
\hline
\end{tabular}
\end{table}



\paragraph{T-EVADE-002: Content Wrapper Escape}



\begin{table}[h]
\centering
\small
\begin{tabular}{|l|l|}
\hline
Attribute & Value \\
\hline
\textbf{ATLAS ID} & AML.T0043 - Craft Adversarial Data \\
\hline
\textbf{Description} & Attacker crafts content that escapes XML wrapper context \\
\hline
\textbf{Attack Vector} & Tag manipulation, context confusion, instruction override \\
\hline
\textbf{Affected Components} & External content wrapping \\
\hline
\textbf{Current Mitigations} & XML tags + security notice \\
\hline
\textbf{Residual Risk} & Medium - Novel escapes discovered regularly \\
\hline
\textbf{Recommendations} & Multiple wrapper layers, output-side validation \\
\hline
\end{tabular}
\end{table}



\hrulefill


\subsubsection{3.6 Discovery (AML.TA0008)}


\paragraph{T-DISC-001: Tool Enumeration}



\begin{table}[h]
\centering
\small
\begin{tabular}{|l|l|}
\hline
Attribute & Value \\
\hline
\textbf{ATLAS ID} & AML.T0040 - AI Model Inference API Access \\
\hline
\textbf{Description} & Attacker enumerates available tools through prompting \\
\hline
\textbf{Attack Vector} & "What tools do you have?" style queries \\
\hline
\textbf{Affected Components} & Agent tool registry \\
\hline
\textbf{Current Mitigations} & None specific \\
\hline
\textbf{Residual Risk} & Low - Tools generally documented \\
\hline
\textbf{Recommendations} & Consider tool visibility controls \\
\hline
\end{tabular}
\end{table}



\paragraph{T-DISC-002: Session Data Extraction}



\begin{table}[h]
\centering
\small
\begin{tabular}{|l|l|}
\hline
Attribute & Value \\
\hline
\textbf{ATLAS ID} & AML.T0040 - AI Model Inference API Access \\
\hline
\textbf{Description} & Attacker extracts sensitive data from session context \\
\hline
\textbf{Attack Vector} & "What did we discuss?" queries, context probing \\
\hline
\textbf{Affected Components} & Session transcripts, context window \\
\hline
\textbf{Current Mitigations} & Session isolation per sender \\
\hline
\textbf{Residual Risk} & Medium - Within-session data accessible \\
\hline
\textbf{Recommendations} & Implement sensitive data redaction in context \\
\hline
\end{tabular}
\end{table}



\hrulefill


\subsubsection{3.7 Collection \& Exfiltration (AML.TA0009, AML.TA0010)}


\paragraph{T-EXFIL-001: Data Theft via web\_fetch}



\begin{table}[h]
\centering
\small
\begin{tabular}{|l|l|}
\hline
Attribute & Value \\
\hline
\textbf{ATLAS ID} & AML.T0009 - Collection \\
\hline
\textbf{Description} & Attacker exfiltrates data by instructing agent to send to external URL \\
\hline
\textbf{Attack Vector} & Prompt injection causing agent to POST data to attacker server \\
\hline
\textbf{Affected Components} & web\_fetch tool \\
\hline
\textbf{Current Mitigations} & SSRF blocking for internal networks \\
\hline
\textbf{Residual Risk} & High - External URLs permitted \\
\hline
\textbf{Recommendations} & Implement URL allowlisting, data classification awareness \\
\hline
\end{tabular}
\end{table}



\paragraph{T-EXFIL-002: Unauthorized Message Sending}



\begin{table}[h]
\centering
\small
\begin{tabular}{|l|l|}
\hline
Attribute & Value \\
\hline
\textbf{ATLAS ID} & AML.T0009 - Collection \\
\hline
\textbf{Description} & Attacker causes agent to send messages containing sensitive data \\
\hline
\textbf{Attack Vector} & Prompt injection causing agent to message attacker \\
\hline
\textbf{Affected Components} & Message tool, channel integrations \\
\hline
\textbf{Current Mitigations} & Outbound messaging gating \\
\hline
\textbf{Residual Risk} & Medium - Gating may be bypassed \\
\hline
\textbf{Recommendations} & Require explicit confirmation for new recipients \\
\hline
\end{tabular}
\end{table}



\paragraph{T-EXFIL-003: Credential Harvesting}



\begin{table}[h]
\centering
\small
\begin{tabular}{|l|l|}
\hline
Attribute & Value \\
\hline
\textbf{ATLAS ID} & AML.T0009 - Collection \\
\hline
\textbf{Description} & Malicious skill harvests credentials from agent context \\
\hline
\textbf{Attack Vector} & Skill code reads environment variables, config files \\
\hline
\textbf{Affected Components} & Skill execution environment \\
\hline
\textbf{Current Mitigations} & None specific to skills \\
\hline
\textbf{Residual Risk} & Critical - Skills run with agent privileges \\
\hline
\textbf{Recommendations} & Skill sandboxing, credential isolation \\
\hline
\end{tabular}
\end{table}



\hrulefill


\subsubsection{3.8 Impact (AML.TA0011)}


\paragraph{T-IMPACT-001: Unauthorized Command Execution}



\begin{table}[h]
\centering
\small
\begin{tabular}{|l|l|}
\hline
Attribute & Value \\
\hline
\textbf{ATLAS ID} & AML.T0031 - Erode AI Model Integrity \\
\hline
\textbf{Description} & Attacker executes arbitrary commands on user system \\
\hline
\textbf{Attack Vector} & Prompt injection combined with exec approval bypass \\
\hline
\textbf{Affected Components} & Bash tool, command execution \\
\hline
\textbf{Current Mitigations} & Exec approvals, Docker sandbox option \\
\hline
\textbf{Residual Risk} & Critical - Host execution without sandbox \\
\hline
\textbf{Recommendations} & Default to sandbox, improve approval UX \\
\hline
\end{tabular}
\end{table}



\paragraph{T-IMPACT-002: Resource Exhaustion (DoS)}



\begin{table}[h]
\centering
\small
\begin{tabular}{|l|l|}
\hline
Attribute & Value \\
\hline
\textbf{ATLAS ID} & AML.T0031 - Erode AI Model Integrity \\
\hline
\textbf{Description} & Attacker exhausts API credits or compute resources \\
\hline
\textbf{Attack Vector} & Automated message flooding, expensive tool calls \\
\hline
\textbf{Affected Components} & Gateway, agent sessions, API provider \\
\hline
\textbf{Current Mitigations} & None \\
\hline
\textbf{Residual Risk} & High - No rate limiting \\
\hline
\textbf{Recommendations} & Implement per-sender rate limits, cost budgets \\
\hline
\end{tabular}
\end{table}



\paragraph{T-IMPACT-003: Reputation Damage}



\begin{table}[h]
\centering
\small
\begin{tabular}{|l|l|}
\hline
Attribute & Value \\
\hline
\textbf{ATLAS ID} & AML.T0031 - Erode AI Model Integrity \\
\hline
\textbf{Description} & Attacker causes agent to send harmful/offensive content \\
\hline
\textbf{Attack Vector} & Prompt injection causing inappropriate responses \\
\hline
\textbf{Affected Components} & Output generation, channel messaging \\
\hline
\textbf{Current Mitigations} & LLM provider content policies \\
\hline
\textbf{Residual Risk} & Medium - Provider filters imperfect \\
\hline
\textbf{Recommendations} & Output filtering layer, user controls \\
\hline
\end{tabular}
\end{table}



\hrulefill


\subsection{4. ClawHub Supply Chain Analysis}


\subsubsection{4.1 Current Security Controls}



\begin{table}[h]
\centering
\small
\begin{tabular}{|l|l|l|}
\hline
Control & Implementation & Effectiveness \\
\hline
GitHub Account Age & \texttt{requireGitHubAccountAge()} & Medium - Raises bar for new attackers \\
\hline
Path Sanitization & \texttt{sanitizePath()} & High - Prevents path traversal \\
\hline
File Type Validation & \texttt{isTextFile()} & Medium - Only text files, but can still be malicious \\
\hline
Size Limits & 50MB total bundle & High - Prevents resource exhaustion \\
\hline
Required SKILL.md & Mandatory readme & Low security value - Informational only \\
\hline
Pattern Moderation & FLAG\_RULES in moderation.ts & Low - Easily bypassed \\
\hline
Moderation Status & \texttt{moderationStatus} field & Medium - Manual review possible \\
\hline
\end{tabular}
\end{table}



\subsubsection{4.2 Moderation Flag Patterns}


Current patterns in \texttt{moderation.ts}:

\begin{lstlisting}[language=javascript]
// Known-bad identifiers
/(keepcold131\/ClawdAuthenticatorTool|ClawdAuthenticatorTool)/i

// Suspicious keywords
/(malware|stealer|phish|phishing|keylogger)/i
/(api[-_ ]?key|token|password|private key|secret)/i
/(wallet|seed phrase|mnemonic|crypto)/i
/(discord\.gg|webhook|hooks\.slack)/i
/(curl[^\n]+\|\s*(sh|bash))/i
/(bit\.ly|tinyurl\.com|t\.co|goo\.gl|is\.gd)/i
\end{lstlisting}


\textbf{Limitations:}

\begin{itemize}
  \item Only checks slug, displayName, summary, frontmatter, metadata, file paths
  \item Does not analyze actual skill code content
  \item Simple regex easily bypassed with obfuscation
  \item No behavioral analysis
\end{itemize}


\subsubsection{4.3 Planned Improvements}



\begin{table}[h]
\centering
\small
\begin{tabular}{|l|l|l|}
\hline
Improvement & Status & Impact \\
\hline
VirusTotal Integration & In Progress & High - Code Insight behavioral analysis \\
\hline
Community Reporting & Partial (\texttt{skillReports} table exists) & Medium \\
\hline
Audit Logging & Partial (\texttt{auditLogs} table exists) & Medium \\
\hline
Badge System & Implemented & Medium - \texttt{highlighted}, \texttt{official}, \texttt{deprecated}, \texttt{redactionApproved} \\
\hline
\end{tabular}
\end{table}



\hrulefill


\subsection{5. Risk Matrix}


\subsubsection{5.1 Likelihood vs Impact}



\begin{table}[h]
\centering
\small
\begin{tabular}{|l|l|l|l|l|}
\hline
Threat ID & Likelihood & Impact & Risk Level & Priority \\
\hline
T-EXEC-001 & High & Critical & \textbf{Critical} & P0 \\
\hline
T-PERSIST-001 & High & Critical & \textbf{Critical} & P0 \\
\hline
T-EXFIL-003 & Medium & Critical & \textbf{Critical} & P0 \\
\hline
T-IMPACT-001 & Medium & Critical & \textbf{High} & P1 \\
\hline
T-EXEC-002 & High & High & \textbf{High} & P1 \\
\hline
T-EXEC-004 & Medium & High & \textbf{High} & P1 \\
\hline
T-ACCESS-003 & Medium & High & \textbf{High} & P1 \\
\hline
T-EXFIL-001 & Medium & High & \textbf{High} & P1 \\
\hline
T-IMPACT-002 & High & Medium & \textbf{High} & P1 \\
\hline
T-EVADE-001 & High & Medium & \textbf{Medium} & P2 \\
\hline
T-ACCESS-001 & Low & High & \textbf{Medium} & P2 \\
\hline
T-ACCESS-002 & Low & High & \textbf{Medium} & P2 \\
\hline
T-PERSIST-002 & Low & High & \textbf{Medium} & P2 \\
\hline
\end{tabular}
\end{table}



\subsubsection{5.2 Critical Path Attack Chains}


\textbf{Attack Chain 1: Skill-Based Data Theft}

\begin{lstlisting}
T-PERSIST-001 → T-EVADE-001 → T-EXFIL-003
(Publish malicious skill) → (Evade moderation) → (Harvest credentials)
\end{lstlisting}


\textbf{Attack Chain 2: Prompt Injection to RCE}

\begin{lstlisting}
T-EXEC-001 → T-EXEC-004 → T-IMPACT-001
(Inject prompt) → (Bypass exec approval) → (Execute commands)
\end{lstlisting}


\textbf{Attack Chain 3: Indirect Injection via Fetched Content}

\begin{lstlisting}
T-EXEC-002 → T-EXFIL-001 → External exfiltration
(Poison URL content) → (Agent fetches & follows instructions) → (Data sent to attacker)
\end{lstlisting}


\hrulefill


\subsection{6. Recommendations Summary}


\subsubsection{6.1 Immediate (P0)}



\begin{table}[h]
\centering
\small
\begin{tabular}{|l|l|l|}
\hline
ID & Recommendation & Addresses \\
\hline
R-001 & Complete VirusTotal integration & T-PERSIST-001, T-EVADE-001 \\
\hline
R-002 & Implement skill sandboxing & T-PERSIST-001, T-EXFIL-003 \\
\hline
R-003 & Add output validation for sensitive actions & T-EXEC-001, T-EXEC-002 \\
\hline
\end{tabular}
\end{table}



\subsubsection{6.2 Short-term (P1)}



\begin{table}[h]
\centering
\small
\begin{tabular}{|l|l|l|}
\hline
ID & Recommendation & Addresses \\
\hline
R-004 & Implement rate limiting & T-IMPACT-002 \\
\hline
R-005 & Add token encryption at rest & T-ACCESS-003 \\
\hline
R-006 & Improve exec approval UX and validation & T-EXEC-004 \\
\hline
R-007 & Implement URL allowlisting for web\_fetch & T-EXFIL-001 \\
\hline
\end{tabular}
\end{table}



\subsubsection{6.3 Medium-term (P2)}



\begin{table}[h]
\centering
\small
\begin{tabular}{|l|l|l|}
\hline
ID & Recommendation & Addresses \\
\hline
R-008 & Add cryptographic channel verification where possible & T-ACCESS-002 \\
\hline
R-009 & Implement config integrity verification & T-PERSIST-003 \\
\hline
R-010 & Add update signing and version pinning & T-PERSIST-002 \\
\hline
\end{tabular}
\end{table}



\hrulefill


\subsection{7. Appendices}


\subsubsection{7.1 ATLAS Technique Mapping}



\begin{table}[h]
\centering
\small
\begin{tabular}{|l|l|l|}
\hline
ATLAS ID & Technique Name & OpenClaw Threats \\
\hline
AML.T0006 & Active Scanning & T-RECON-001, T-RECON-002 \\
\hline
AML.T0009 & Collection & T-EXFIL-001, T-EXFIL-002, T-EXFIL-003 \\
\hline
AML.T0010.001 & Supply Chain: AI Software & T-PERSIST-001, T-PERSIST-002 \\
\hline
AML.T0010.002 & Supply Chain: Data & T-PERSIST-003 \\
\hline
AML.T0031 & Erode AI Model Integrity & T-IMPACT-001, T-IMPACT-002, T-IMPACT-003 \\
\hline
AML.T0040 & AI Model Inference API Access & T-ACCESS-001, T-ACCESS-002, T-ACCESS-003, T-DISC-001, T-DISC-002 \\
\hline
AML.T0043 & Craft Adversarial Data & T-EXEC-004, T-EVADE-001, T-EVADE-002 \\
\hline
AML.T0051.000 & LLM Prompt Injection: Direct & T-EXEC-001, T-EXEC-003 \\
\hline
AML.T0051.001 & LLM Prompt Injection: Indirect & T-EXEC-002 \\
\hline
\end{tabular}
\end{table}



\subsubsection{7.2 Key Security Files}



\begin{table}[h]
\centering
\small
\begin{tabular}{|l|l|l|}
\hline
Path & Purpose & Risk Level \\
\hline
\texttt{src/infra/exec-approvals.ts} & Command approval logic & \textbf{Critical} \\
\hline
\texttt{src/gateway/auth.ts} & Gateway authentication & \textbf{Critical} \\
\hline
\texttt{src/web/inbound/access-control.ts} & Channel access control & \textbf{Critical} \\
\hline
\texttt{src/infra/net/ssrf.ts} & SSRF protection & \textbf{Critical} \\
\hline
\texttt{src/security/external-content.ts} & Prompt injection mitigation & \textbf{Critical} \\
\hline
\texttt{src/agents/sandbox/tool-policy.ts} & Tool policy enforcement & \textbf{Critical} \\
\hline
\texttt{convex/lib/moderation.ts} & ClawHub moderation & \textbf{High} \\
\hline
\texttt{convex/lib/skillPublish.ts} & Skill publishing flow & \textbf{High} \\
\hline
\texttt{src/routing/resolve-route.ts} & Session isolation & \textbf{Medium} \\
\hline
\end{tabular}
\end{table}



\subsubsection{7.3 Glossary}



\begin{table}[h]
\centering
\small
\begin{tabular}{|l|l|}
\hline
Term & Definition \\
\hline
\textbf{ATLAS} & MITRE's Adversarial Threat Landscape for AI Systems \\
\hline
\textbf{ClawHub} & OpenClaw's skill marketplace \\
\hline
\textbf{Gateway} & OpenClaw's message routing and authentication layer \\
\hline
\textbf{MCP} & Model Context Protocol - tool provider interface \\
\hline
\textbf{Prompt Injection} & Attack where malicious instructions are embedded in input \\
\hline
\textbf{Skill} & Downloadable extension for OpenClaw agents \\
\hline
\textbf{SSRF} & Server-Side Request Forgery \\
\hline
\end{tabular}
\end{table}



\hrulefill


\_This threat model is a living document. Report security issues to security@openclaw.ai\_


\clearpage

% === Source file: security/formal-verification.md ===
\section{Formal Verification (Security Models)}


This page tracks OpenClaw’s \textbf{formal security models} (TLA+/TLC today; more as needed).

\begin{quote}
Note: some older links may refer to the previous project name.
\end{quote}


\textbf{Goal (north star):} provide a machine-checked argument that OpenClaw enforces its
intended security policy (authorization, session isolation, tool gating, and
misconfiguration safety), under explicit assumptions.

\textbf{What this is (today):} an executable, attacker-driven \textbf{security regression suite}:

\begin{itemize}
  \item Each claim has a runnable model-check over a finite state space.
  \item Many claims have a paired \textbf{negative model} that produces a counterexample trace for a realistic bug class.
\end{itemize}


\textbf{What this is not (yet):} a proof that “OpenClaw is secure in all respects” or that the full TypeScript implementation is correct.

\subsection{Where the models live}


Models are maintained in a separate repo: \href{https://github.com/vignesh07/openclaw-formal-models}{vignesh07/openclaw-formal-models}.

\subsection{Important caveats}


\begin{itemize}
  \item These are \textbf{models}, not the full TypeScript implementation. Drift between model and code is possible.
  \item Results are bounded by the state space explored by TLC; “green” does not imply security beyond the modeled assumptions and bounds.
  \item Some claims rely on explicit environmental assumptions (e.g., correct deployment, correct configuration inputs).
\end{itemize}


\subsection{Reproducing results}


Today, results are reproduced by cloning the models repo locally and running TLC (see below). A future iteration could offer:

\begin{itemize}
  \item CI-run models with public artifacts (counterexample traces, run logs)
  \item a hosted “run this model” workflow for small, bounded checks
\end{itemize}


Getting started:

\begin{lstlisting}[language=bash]
git clone https://github.com/vignesh07/openclaw-formal-models
cd openclaw-formal-models

# Java 11+ required (TLC runs on the JVM).
# The repo vendors a pinned `tla2tools.jar` (TLA+ tools) and provides `bin/tlc` + Make targets.

make <target>
\end{lstlisting}


\subsubsection{Gateway exposure and open gateway misconfiguration}


\textbf{Claim:} binding beyond loopback without auth can make remote compromise possible / increases exposure; token/password blocks unauth attackers (per the model assumptions).

\begin{itemize}
  \item Green runs:
  \item \texttt{make gateway-exposure-v2}
  \item \texttt{make gateway-exposure-v2-protected}
  \item Red (expected):
  \item \texttt{make gateway-exposure-v2-negative}
\end{itemize}


See also: \texttt{docs/gateway-exposure-matrix.md} in the models repo.

\subsubsection{Nodes.run pipeline (highest-risk capability)}


\textbf{Claim:} \texttt{nodes.run} requires (a) node command allowlist plus declared commands and (b) live approval when configured; approvals are tokenized to prevent replay (in the model).

\begin{itemize}
  \item Green runs:
  \item \texttt{make nodes-pipeline}
  \item \texttt{make approvals-token}
  \item Red (expected):
  \item \texttt{make nodes-pipeline-negative}
  \item \texttt{make approvals-token-negative}
\end{itemize}


\subsubsection{Pairing store (DM gating)}


\textbf{Claim:} pairing requests respect TTL and pending-request caps.

\begin{itemize}
  \item Green runs:
  \item \texttt{make pairing}
  \item \texttt{make pairing-cap}
  \item Red (expected):
  \item \texttt{make pairing-negative}
  \item \texttt{make pairing-cap-negative}
\end{itemize}


\subsubsection{Ingress gating (mentions + control-command bypass)}


\textbf{Claim:} in group contexts requiring mention, an unauthorized “control command” cannot bypass mention gating.

\begin{itemize}
  \item Green:
  \item \texttt{make ingress-gating}
  \item Red (expected):
  \item \texttt{make ingress-gating-negative}
\end{itemize}


\subsubsection{Routing/session-key isolation}


\textbf{Claim:} DMs from distinct peers do not collapse into the same session unless explicitly linked/configured.

\begin{itemize}
  \item Green:
  \item \texttt{make routing-isolation}
  \item Red (expected):
  \item \texttt{make routing-isolation-negative}
\end{itemize}


\subsection{v1++: additional bounded models (concurrency, retries, trace correctness)}


These are follow-on models that tighten fidelity around real-world failure modes (non-atomic updates, retries, and message fan-out).

\subsubsection{Pairing store concurrency / idempotency}


\textbf{Claim:} a pairing store should enforce \texttt{MaxPending} and idempotency even under interleavings (i.e., “check-then-write” must be atomic / locked; refresh shouldn’t create duplicates).

What it means:

\begin{itemize}
  \item Under concurrent requests, you can’t exceed \texttt{MaxPending} for a channel.
  \item Repeated requests/refreshes for the same \texttt{(channel, sender)} should not create duplicate live pending rows.
\end{itemize}


\begin{itemize}
  \item Green runs:
  \item \texttt{make pairing-race} (atomic/locked cap check)
  \item \texttt{make pairing-idempotency}
  \item \texttt{make pairing-refresh}
  \item \texttt{make pairing-refresh-race}
  \item Red (expected):
  \item \texttt{make pairing-race-negative} (non-atomic begin/commit cap race)
  \item \texttt{make pairing-idempotency-negative}
  \item \texttt{make pairing-refresh-negative}
  \item \texttt{make pairing-refresh-race-negative}
\end{itemize}


\subsubsection{Ingress trace correlation / idempotency}


\textbf{Claim:} ingestion should preserve trace correlation across fan-out and be idempotent under provider retries.

What it means:

\begin{itemize}
  \item When one external event becomes multiple internal messages, every part keeps the same trace/event identity.
  \item Retries do not result in double-processing.
  \item If provider event IDs are missing, dedupe falls back to a safe key (e.g., trace ID) to avoid dropping distinct events.
\end{itemize}


\begin{itemize}
  \item Green:
  \item \texttt{make ingress-trace}
  \item \texttt{make ingress-trace2}
  \item \texttt{make ingress-idempotency}
  \item \texttt{make ingress-dedupe-fallback}
  \item Red (expected):
  \item \texttt{make ingress-trace-negative}
  \item \texttt{make ingress-trace2-negative}
  \item \texttt{make ingress-idempotency-negative}
  \item \texttt{make ingress-dedupe-fallback-negative}
\end{itemize}


\subsubsection{Routing dmScope precedence + identityLinks}


\textbf{Claim:} routing must keep DM sessions isolated by default, and only collapse sessions when explicitly configured (channel precedence + identity links).

What it means:

\begin{itemize}
  \item Channel-specific dmScope overrides must win over global defaults.
  \item identityLinks should collapse only within explicit linked groups, not across unrelated peers.
\end{itemize}


\begin{itemize}
  \item Green:
  \item \texttt{make routing-precedence}
  \item \texttt{make routing-identitylinks}
  \item Red (expected):
  \item \texttt{make routing-precedence-negative}
  \item \texttt{make routing-identitylinks-negative}
\end{itemize}



\clearpage
